% Example document for trying the report viewer: python3 report_viewer.py docs
\documentclass[11pt]{article}
\usepackage[margin=2.5cm]{geometry}
\usepackage{amsmath}
\title{Example Report: Timing a Pendulum}
\author{Report Viewer example}
\date{}
\begin{document}
\maketitle

\section{Purpose}
This short document exists so you can try the report viewer. Select a sentence and add a comment, press
\emph{Read aloud} and tap any word to hear the text from there, or ask a Claude Code session to answer your comments.

\section{Method}
A small mass hangs from a light string of length $L$. For small swings the period is
\begin{equation}
  T = 2\pi\sqrt{\frac{L}{g}},
\end{equation}
where $g$ is the acceleration due to gravity. The read-aloud skips display equations like this one and carries
on with the next sentence. Timing ten swings with a stopwatch and dividing by ten reduces the effect of reaction
time on each measurement.

\section{What to try}
Change the length of the string, time the swings again, and plot the period against the square root of the
length. The points should fall on a straight line through the origin, and its slope gives an estimate of $g$.
Write down the uncertainty of every timing, because the scatter between repeated measurements shows how far
the estimate can be trusted.

\end{document}
